% !TEX root = ../main.tex
\chapter{Preface}

Could we build a shortcut through space? Could a ship cross the galaxy while
its passengers age only a few years, or a tunnel connect two stars? General
relativity gives these questions a precise form. It tells us that space and
time together make up a geometry that matter can bend, and it gives the rule
connecting the two. Once that rule is in hand, the question stops being
\enquote{is this allowed by physics?} in the abstract. It becomes an
engineering question: what geometry would do the job, what would it take to
produce that geometry, and can nature supply what it takes?

This book teaches you to work through that question. By the end of it you
will be able to write down a spacetime metric designed to do something, work out the
matter and energy it requires, test that requirement against everything
physics knows about matter, and judge how far the design is from anything
that could be built. Along the way you will meet the classic designs of the
field: traversable wormholes, warp drives and Krasnikov tubes, together with
the results that constrain them and the questions that remain open.

\section*{Who this book is for}

The book is written for upper-level undergraduates in physics, graduate
students in engineering, and anyone with similar preparation. You should be
comfortable with special relativity at the level of an introductory modern
physics course, with multivariable calculus and linear algebra, and with
solving ordinary differential equations. You do not need to have studied
general relativity: the first chapters develop the parts of it this book
uses, starting from spacetime diagrams. Readers who already know general
relativity can move quickly through those chapters and will find the
material aimed at them in the boxes marked \emph{Going deeper}.

\section*{How the book is organized}

Part~I builds the foundations: what spacetime engineering is, the geometry of
spacetime, the way matter shapes that geometry, and the physical limits on
matter that decide what can be built. Part~II studies the classic designs one
at a time. Part~III turns to design itself: how each part of a spacetime's
geometry affects what it demands and what it does to travellers. Part~IV asks
what kinds of matter and fields could supply those demands. Part~V takes up
dynamics, the checking of calculations, and the open problems of the field.

\section*{How to use this book}

Each chapter opens with a short list of what you will learn and closes with a
summary. Worked examples show each method carried through in full, with the
reasoning spelled out. The \emph{Check your understanding} questions in the
text are short; try each one before reading on, and compare with the answers
at the end of the chapter. The problems at the end of each chapter are marked
by difficulty: [$\star$] problems practise a single idea, [$\star\star$]
problems combine several, and [$\star\star\star$] problems are open-ended or
demanding. Boxes marked \emph{Going deeper} carry material for readers with
more background, and can be skipped on a first reading.
